%==========================================================
% PART 2: the AI model (Phase 2). Added after Phase 2 was run.
% Everything before this point is unchanged.
%==========================================================
% generated by scripts/export_explainer_varied.py -- do not edit
\newcommand{\SeedRows}{0 & $-$12\% & 0.87 & $-$28\% & 0.69 & 10\% / 10\%\\
1 & $+$7\% & 1.01 & $-$14\% & 0.79 & 12\% / 16\%\\
2 & $+$16\% & 1.12 & $-$5\% & 0.89 & 11\% / 16\%\\
3 & $+$6\% & 1.02 & $-$18\% & 0.77 & 11\% / 14\%\\
4 & $+$5\% & 1.02 & $-$14\% & 0.81 & 12\% / 15\%\\}
\newcommand{\SeedValTestR}{0.97}
\newcommand{\ExpRows}{Sweden only (as before) & 0.62 & $-$32\% to $-$6\% & $-$6\% to $+$41\% & 83\% & 32\%\\
Sweden only, rate error in log scale & 0.62 & $-$20\% to $+$9\% & $-$8\% to $+$32\% & 81\% & 48\%\\
Sweden + Italy 2021 & 0.65 & $-$17\% to $+$0\% & $-$3\% to $+$4\% & 95\% & 28\%\\
Sweden + Italy 2021, log scale & 0.65 & $-$23\% to $-$16\% & $-$15\% to $-$5\% & 83\% & 20\%\\
Sweden + Italy 2021 + neighbour inputs & 0.68 & $-$17\% to $-$11\% & $-$16\% to $+$0\% & 91\% & 35\%\\
Sweden + Italy 2021, reliable links only & 0.65 & $-$18\% to $-$13\% & $-$7\% to $+$5\% & 95\% & 42\%\\
\quad + hidden + rate error on 15-minute averages (not kept) & 0.65 & $-$24\% to $-$13\% & $-$8\% to $-$2\% & 91\% & 50\%\\
\quad + Swedish radar corrected with gauges & 0.65 & $-$13\% to $-$9\% & $-$4\% to $+$12\% & 94\% & 40\%\\
\quad + unclear moments hidden & 0.65 & $-$15\% to $-$11\% & $-$5\% to $+$5\% & 93\% & 56\%\\
\quad + both of the above & 0.65 & $-$13\% to $-$6\% & $+$1\% to $+$11\% & 95\% & 55\%\\
\quad + hidden + gauge answers in Italy & 0.65 & $-$15\% to $+$2\% & $+$5\% to $+$11\% & 97\% & 67\%\\
\textbf{\quad + all three (final recipe)} & 0.65 & $-$11\% to $+$1\% & $+$5\% to $+$13\% & 96\% & 66\%\\
\quad + shorter memory (not kept) & 0.66 & $-$14\% to $-$4\% & $-$2\% to $+$5\% & 95\% & 51\%\\
\quad + corrected + hidden + shorter memory (not kept) & 0.65 & $-$8\% to $-$3\% & $-$2\% to $+$8\% & 93\% & 59\%\\}
\newcommand{\ExpRelSeRange}{$-$18\% and $-$13\%}
\newcommand{\ExpRelLong}{42\%}
\newcommand{\ExpFinSeRange}{$-$11\% and $+$1\%}
\newcommand{\ExpFinLong}{66\%}
\newcommand{\NlRelWestLong}{66\%}
\newcommand{\NlRelWestGauges}{91\%}
\newcommand{\NlRelEastLong}{22\%}
\newcommand{\NlRelEastGauges}{84\%}
\newcommand{\NlFinWestLong}{93\%}
\newcommand{\NlFinWestGauges}{93\%}
\newcommand{\NlFinEastLong}{30\%}
\newcommand{\NlFinEastGauges}{87\%}
\newcommand{\NlRows}{reliable links only & west $\to$ east & $+$4\% to $+$19\% & 91\% & 104\% & 66\%\\
reliable links only & east $\to$ west & $-$13\% to $-$12\% & 84\% & 84\% & 22\%\\
\textbf{final recipe} & west $\to$ east & $+$9\% to $+$37\% & 93\% & 103\% & 93\%\\
\textbf{final recipe} & east $\to$ west & $-$13\% to $-$6\% & 87\% & 82\% & 30\%\\}
\newcommand{\CVallSeMCC}{0.651}
\newcommand{\CVallSeR}{0.83}
\newcommand{\CVallSeRMSE}{0.97}
\newcommand{\CVallSeBias}{$-$14\%}
\newcommand{\CVallSeBiasRange}{$-$19\% to $-$6\%}
\newcommand{\CVallBiasA}{$+$3\% to $+$15\%}
\newcommand{\CVallLongAa}{94\%}
\newcommand{\CVallLongAb}{94\%}
\newcommand{\CVallLongAc}{86\%}
\newcommand{\CVallLongAd}{96\%}
\newcommand{\CVallGaugesA}{88\%}
\newcommand{\CVallBiasB}{$-$1\% to $+$8\%}
\newcommand{\CVallLongBa}{99\%}
\newcommand{\CVallLongBb}{102\%}
\newcommand{\CVallLongBc}{86\%}
\newcommand{\CVallLongBd}{50\%}
\newcommand{\CVallGaugesB}{84\%}
\newcommand{\CVrelSeMCC}{0.648}
\newcommand{\CVrelSeR}{0.81}
\newcommand{\CVrelSeRMSE}{1.01}
\newcommand{\CVrelSeBias}{$-$17\%}
\newcommand{\CVrelSeBiasRange}{$-$23\% to $-$12\%}
\newcommand{\CVrelBiasA}{$+$1\% to $+$4\%}
\newcommand{\CVrelLongAa}{93\%}
\newcommand{\CVrelLongAb}{94\%}
\newcommand{\CVrelLongAc}{87\%}
\newcommand{\CVrelLongAd}{95\%}
\newcommand{\SeOnlyLongAa}{77\%}
\newcommand{\SeOnlyLongAb}{79\%}
\newcommand{\SeOnlyLongAc}{70\%}
\newcommand{\SeOnlyLongAd}{77\%}
\newcommand{\CVrelGaugesA}{87\%}
\newcommand{\CVrelBiasB}{$-$1\% to $+$4\%}
\newcommand{\CVrelLongBa}{98\%}
\newcommand{\CVrelLongBb}{102\%}
\newcommand{\CVrelLongBc}{87\%}
\newcommand{\CVrelLongBd}{60\%}
\newcommand{\SeOnlyLongBa}{84\%}
\newcommand{\SeOnlyLongBb}{86\%}
\newcommand{\SeOnlyLongBc}{72\%}
\newcommand{\SeOnlyLongBd}{55\%}
\newcommand{\CVrelGaugesB}{86\%}
\newcommand{\CVfinSeMCC}{0.645}
\newcommand{\CVfinSeR}{0.85}
\newcommand{\CVfinSeRMSE}{0.91}
\newcommand{\CVfinSeBias}{$-$10\%}
\newcommand{\CVfinSeBiasRange}{$-$13\% to $-$5\%}
\newcommand{\CVfinBiasA}{$+$5\% to $+$13\%}
\newcommand{\CVfinLongAa}{93\%}
\newcommand{\CVfinLongAb}{98\%}
\newcommand{\CVfinLongAc}{97\%}
\newcommand{\CVfinLongAd}{107\%}
\newcommand{\CVfinGaugesA}{90\%}
\newcommand{\CVfinBiasB}{$+$6\% to $+$12\%}
\newcommand{\CVfinLongBa}{98\%}
\newcommand{\CVfinLongBb}{107\%}
\newcommand{\CVfinLongBc}{95\%}
\newcommand{\CVfinLongBd}{74\%}
\newcommand{\CVfinGaugesB}{89\%}
\newcommand{\Mayradar}{51}
\newcommand{\Maygauge}{48}
\newcommand{\Mayaise}{12}
\newcommand{\Mayaicv}{17}
\newcommand{\Mayaifin}{25}

\clearpage
\begin{center}
\begin{minipage}{0.86\textwidth}
\small
\textbf{Part 2: the AI, and what happened when we tested it.} The first part of this document explained how a
radio link feels rain and how the classical method turns that into millimetres per hour. This second part explains
the neural network we built to do the same job, from the very first idea of what a neural network is, then shows
every result: in Sweden, where it was trained, and in Italy and the Netherlands, which it had never seen. As before,
every word is explained, and the numbers are real.
\end{minipage}
\end{center}

%==========================================================
\section{What a neural network is, from scratch}
\label{sec:nn}
%==========================================================

In Section~\ref{sec:inverse}, people wrote the rules: ``call it rain if the signal's jumpiness is above 0.35\,dB'',
``freeze the baseline at the average of the last 5 dry minutes''. A \textbf{neural network} is a different way to get
rules. Instead of writing them, we show the computer many examples of the input (the signal) together with the right
answer (what the radar measured at the same time), and let it adjust itself until its answers match. The rules it ends
up with are not written in words; they are stored as a long list of numbers.

\subsection{The smallest piece: one ``neuron''}
A neuron is a tiny calculator. It takes a few numbers in, multiplies each by its own \textbf{weight}, adds them up,
adds one more number called the \textbf{bias}, and then \emph{bends} the result. Bending means: negative results are
squashed towards zero, positive ones pass through. (Our network uses a smooth version of this bend called GELU; the
idea is the same.)

\begin{example}[: one neuron]
Inputs: the link lost 9.2\,dB more than usual, and the signal's jumpiness is 1.5\,dB. Weights $0.4$ and $1.0$, bias $-3$.
\[
0.4\times 9.2 \;+\; 1.0\times 1.5 \;+\; (-3) \;=\; 3.68 + 1.5 - 3 \;=\; 2.18 .
\]
The result is positive, so it passes the bend almost unchanged: the neuron ``fires'' with strength about 2.2.
On a dry day (0.1\,dB extra loss, jumpiness 0.1\,dB) the same neuron gives $0.04+0.1-3=-2.86$, which the bend squashes
to nearly zero: it stays quiet. This one neuron has, in effect, become a little rain detector.
\end{example}

A network is many neurons arranged in \textbf{layers}: the outputs of one layer are the inputs of the next. Early
layers notice simple things (``the signal dropped''), later layers combine them into more complex things (``it dropped,
stayed down for twenty minutes, and wobbled the way rain makes it wobble''). Our network has 111{,}890 weights and
biases in total. We will call them its \emph{dials}.

\subsection{Turning a number into a probability}
To answer ``is it raining?'' we want a number between 0 and 1, a probability. The last step does this with the
\textbf{sigmoid} function,
\[
p = \frac{1}{1+e^{-x}} ,
\]
which turns any number $x$ into something between 0 and 1: large negative $x$ gives almost 0, $x=0$ gives exactly 0.5,
large positive $x$ gives almost 1. ($e \approx 2.718$ is a fixed mathematical constant.) For $x=2.18$:
$p = 1/(1+e^{-2.18}) = 1/(1+0.113) = 0.90$, ``90\% sure it is raining''.

\subsection{Learning: turning the dials}
At first the dials are set to random values and the network's answers are nonsense. Learning goes in rounds:
\begin{enumerate}
\item Give the network a handful of examples and let it answer.
\item Measure how wrong it was with one number, the \textbf{loss} (here this word means ``penalty for being wrong'',
not signal loss in dB).
\item For every dial, work out whether turning it slightly up or slightly down would make the penalty smaller.
The mathematics for this (called \emph{back-propagation}) does this for all 111{,}890 dials at once.
\item Turn every dial a tiny bit in its helpful direction. Repeat.
\end{enumerate}
After thousands of rounds the dials settle at values that give small penalties on the examples.

\begin{example}[: the two penalties we use]
\textbf{For rain/no-rain}, the penalty is
\[ -\ln(\text{the probability the network gave to what really happened}), \]
where $\ln$ is the natural logarithm. If it said 90\% rain and it rained: $-\ln 0.9 = 0.105$, a small penalty.
If it said 10\% rain and it rained: $-\ln 0.1 = 2.30$, a big penalty. Being confidently wrong hurts most.

\textbf{For the amount}, the penalty is the squared mistake. Radar says 8\,mm/h, network says 5: penalty
$(8-5)^2 = 9$. Network says 7: penalty 1. Squaring makes big misses count much more than small ones, the same idea as the
RMSE in Section~\ref{sec:scores}.
\end{example}

\subsection{Three piles of data, like school}
If a student sees the exam questions while studying, a perfect exam score proves nothing. The same is true for
networks, so we split time into three piles, exactly as in Phase~1:
\begin{itemize}
\item \textbf{Training} (1 June -- 15 July 2015): the examples the dials are turned on. The homework.
\item \textbf{Validation} (16 -- 31 July): never used to turn dials, only to decide when to stop and which version to
keep. The practice exam.
\item \textbf{Test} (August): scored once at the end. The real exam.
\end{itemize}
A network that does very well on training data but badly on new data has \textbf{overfitted}: it memorised the homework
instead of learning the subject. Watching the validation score is how we notice.

\begin{oneline}
A neural network is a big adjustable formula; training means turning its dials until its answers on known examples are
good, and we only trust it on examples it never saw.
\end{oneline}

%==========================================================
\section{Our network: a reader that only looks backwards}
\label{sec:tcn}
%==========================================================

\subsection{What it reads}
Every minute, for every link, the network is given five numbers about the signal and three facts about the link.

\begin{center}\small
\begin{tabular}{p{3.6cm}p{10.4cm}}
\toprule
Input & What it is, in words \\
\midrule
Extra loss & The total loss this minute minus the \textbf{dry level}: a low value of the total loss over the previous
3 days (explained in Section~\ref{sec:drylevel} below). \\
Physics hint & The rain rate that this extra loss would mean if all of it were rain, from Equation~\ref{eq:inverse}. \\
Jumpiness & The standard deviation of the total loss over the previous 15 minutes (Step~2 of Section~\ref{sec:inverse}). \\
Present & 1 if the link reported this minute, 0 if the value is missing. \\
Extra loss, 24-hour view & The total loss minus the \emph{median} of the previous 24 hours. Next to the first input, it tells
the network whether the last day was mostly rainy (Section~\ref{sec:drylevel}). \\
\midrule
Link facts & Its length, its frequency, its polarisation. These never change. \\
\bottomrule
\end{tabular}
\end{center}

Notice the word \emph{previous} everywhere. Nothing the network reads comes from the future. The classical method of
Part~1 needed 30 minutes of future signal to decide whether it rained; ours does not, so it can run live.

\subsection{Finding the dry level, even after a rainy day}
\label{sec:drylevel}
Everything rests on knowing what the link's total loss would be \emph{if it were dry right now}. The classical method
of Part~1 freezes the level of the last dry minutes before a shower. The network needs something that works every
minute, with no wet/dry decision made beforehand.

\textbf{First attempt: the 24-hour median.} Sort the total loss of the previous 24 hours and take the middle value. A
storm of a few hours barely moves it. But there is a trap. If it rained for \emph{more than half} of those 24 hours, the
middle value is itself a rainy value. The dry level then sits too high, the extra loss looks too small, and the method
under-measures, exactly when the ground is already soaked and floods are most likely.

We checked how often this happens: in the Italian summer of 2022 and the Dutch summer of 2014, about one rainy hour in
ten came after a day that had been mostly rainy (in Sweden in August, almost never). In those hours the network with
the 24-hour median measured only 48\% of the Dutch rain, against about 100\% at other times.

\textbf{The fix: a low value over 3 days.} Rain only ever \emph{adds} loss; it never makes the beam stronger. So
instead of the middle value we take a \emph{low} one: cut the past 3 days into 15-minute blocks, take each block's
median (to calm the noise), sort the 288 blocks, and take the value 20\% of the way up from the bottom. This is called
the \textbf{20th percentile}. It stays at the dry level as long as at least a fifth of the 3 days was dry: it would
take rain in more than 80\% of three whole days to fool it. We chose 3 days and 20\% from this reasoning, not by trying
values on the test data.

\begin{figure}[h]
\centering
\begin{tikzpicture}[x=0.42cm,y=0.42cm]
  \draw[inkc!40] (0,0) -- (30,0);
  \draw[-{Stealth},inkc!60] (0,-0.8) -- (30.5,-0.8) node[right,font=\scriptsize] {3 days};
  % dry stretches at 1, rainy stretches at 3-6
  \draw[rainc,thick] (0,1) -- (4,1) -- (4.4,4.5) -- (9,4.2) -- (9.4,1.1) -- (11,1) -- (11.4,5.2) -- (19,4.6)
     -- (19.4,1) -- (21,1.1) -- (21.4,3.8) -- (30,4.4);
  \draw[gaugec,thick,dashed] (0,1.05) -- (30,1.05) node[right,font=\scriptsize,align=left] {low value\\(20th percentile)};
  \draw[inkc,thick,dotted] (0,4.2) -- (30,4.2) node[right,font=\scriptsize,align=left] {middle value\\(median)};
  \node[font=\scriptsize,rainc,anchor=south] at (15,5.3) {rain};
  \node[font=\scriptsize,rainc,anchor=north] at (10.2,0.9) {dry};
  \node[font=\scriptsize,rotate=90,anchor=south] at (-0.2,2.8) {total loss};
\end{tikzpicture}
\caption{Illustration (not real data): a link during three mostly rainy days. Because rain filled more than half of the
time, the middle value (dotted) lies on the rainy level. The low value (dashed) still finds the dry level.}
\label{fig:drylevel}
\end{figure}

The old 24-hour median is not thrown away: it becomes the fifth input. When the two disagree, the network can tell that
the recent past was rainy.

\begin{example}[: the extra loss at 08:30 on 25 August]
On sub-link 718 the 3-day low value before 08:30 was 48.70\,dB (the 24-hour median was 48.98). At 08:30 the total
loss was 58.18\,dB, so the extra loss is $58.18-48.70 = 9.48$\,dB. The classical baseline of
Section~\ref{sec:inverse} gave 9.47\,dB: on this ordinary day, the two ways of finding the dry level agree to within
0.01\,dB.
\end{example}

\textbf{Did it work?} Here is how much of the reference rain the AI measured in the hours that followed a mostly
rainy day (more than half of the previous 24 hours rainy). For comparison, the last column shows the same AI at other
times (after a day with less than 10\% rain). The test countries are described in Sections~\ref{sec:italy}
and~\ref{sec:nl}.
\begin{center}\small\setlength{\tabcolsep}{4pt}
\begin{tabular}{lccccc}
\toprule
& share of rainy hours & \multicolumn{3}{c}{AI, after a mostly rainy day} & AI, after a\\
\cmidrule(lr){3-5}
& that follow & & 3-day & + varied & mostly dry\\
Test & a rainy day & 24-h median & low value & training$^{*}$ & day\\
\midrule
Netherlands 2014 & 10\% & 48\% & \textbf{69\%} & -- & 100\%\\
Italy 2022 & 11\% & 64\% & 55\% & \CVfinLongBd & 83\%\\
Italy 2021 & 2\% & 86\% & 77\% & \textbf{\CVfinLongAd} & 77\%\\
\bottomrule
\end{tabular}

\smallskip
{\footnotesize $^{*}$ Section~\ref{sec:varied}: trained on Sweden and on other months of Italy, never on the month judged, with cleaned training answers (Section~\ref{sec:recipe}).}
\end{center}
In the Netherlands the fix works: the shortfall after a rainy day shrank from about half to about a third. In Italy it
did not help: there the AI measures between half and two-thirds of the rain after a rainy day, with either dry level, and the
classical method, which does not use a 24-hour median at all, does no better (62\%). Section~\ref{sec:varied} finds out
why. In short: the dry level is fine; in a long rain the antennas stay wet, the radar we learn from misses the light rain at
its end, and faulty links in the training data had taught the AI that a loss which stays raised for hours is not rain.
Training on more varied weather with cleaned answers (fifth column) fixes this in 2021 and helps clearly in 2022, where these hours come almost
all from one long storm. (An earlier check here said this rain is not lighter than usual. Looked at in 15-minute steps, it is: 1.2
against 2.3\,mm/h on average in 2022, which makes it harder to see under the wet antennas.)

\subsection{How it reads: a magnifying glass sliding along time}
The network's layers are \textbf{convolutions}. A convolution is a small set of weights that slides along the
signal, looking at a few neighbouring minutes at a time, and produces one new number per minute. Picture a magnifying
glass that shows three minutes at once, moved one minute at a time.

Three minutes is short. To see longer patterns without needing a huge number of dials, each layer spaces its three
minutes further apart: the first layer looks at minutes next to each other, the second at minutes 2 apart, then 4, 8,
16, 32, 64 and 128 apart. This trick is called \textbf{dilation}. Stacking eight such layers (each used twice) lets the
final answer depend on the last \textbf{1{,}021 minutes}, about 17 hours. That span is called the
\textbf{receptive field}. Every glass sits at or behind the current minute, never ahead of it; a network built this way
is called a \emph{causal temporal convolutional network} (TCN).

\begin{figure}[h]
\centering
\begin{tikzpicture}[x=0.6cm,y=1.05cm]
  \foreach \row/\lab in {0/{signal (one dot per minute)},1/{layer 1: neighbours},2/{layer 2: 2 apart},3/{layer 3: 4 apart}} {
    \foreach \i in {0,...,15} { \fill[inkc!25] (\i,\row) circle (2.2pt); }
    \node[font=\scriptsize,anchor=east,inkc!80] at (-0.5,\row) {\lab};
  }
  % connections into the last output
  \foreach \s in {15,11,7,3} { \draw[beam!70,thick] (\s,2) -- (15,3); }
  \foreach \t in {15,11,7,3} { \foreach \s in {0,2} { \pgfmathtruncatemacro{\a}{\t-\s} \draw[beam!45] (\a,1) -- (\t,2); } }
  \foreach \t in {15,13,11,9,7,5,3,1} { \foreach \s in {0,1} { \pgfmathtruncatemacro{\a}{\t-\s} \ifnum\a>-1 \draw[beam!30] (\a,0) -- (\t,1); \fi } }
  \fill[gaugec] (15,3) circle (3.2pt);
  \node[font=\scriptsize,anchor=west,gaugec] at (15.3,3) {answer for ``now''};
  \draw[-{Stealth},inkc!60] (0,-0.55) -- (15,-0.55) node[right,font=\scriptsize] {time};
  \node[font=\scriptsize,inkc!60] at (15,-0.85) {now};
\end{tikzpicture}
\caption{How dilation widens the view (simplified: two taps per glass, three layers). Each layer looks further apart, so
the answer for ``now'' (orange) depends on many past minutes but never on future ones. Our real network uses three taps
and eight levels, reaching 1{,}021 minutes into the past.}
\label{fig:tcn}
\end{figure}

\subsection{What it answers: two numbers, then physics}
For each minute the network gives two numbers:
\begin{enumerate}
\item \textbf{The chance it is raining}, $p$, between 0 and 1 (through the sigmoid above). If $p$ is above a
\textbf{threshold} we say ``raining''. We tried thresholds from 0.05 to 0.80 on the validation fortnight and kept the
one with the best MCC: \textbf{0.25}.
\item \textbf{A correction} $c$ in decibels: how much of the extra loss is \emph{not} rain in the air (wet antennas,
the dry level drifting, slow wobbles). It is built to stay between $-1$ and $+4$\,dB, roughly the size wet antennas
really reach.
\end{enumerate}
Then the network steps aside and the physics of Part~1 takes over: the rain-caused loss is $A = \text{extra loss} - c$,
and Equation~\ref{eq:inverse} turns it into mm/h with the ITU coefficients $a$ and $b$ of that link. If $p$ is below the
threshold, the answer is 0.
\[
R = \begin{cases} \left(\dfrac{\text{extra loss} - c}{a\,L}\right)^{1/b} & \text{if } p > 0.25,\\[2ex] 0 & \text{otherwise.}\end{cases}
\]

\begin{oneline}
The AI decides whether it is raining and how much of the loss to blame on things other than rain; the century-old
physics of raindrops still does the conversion into millimetres per hour.
\end{oneline}

\begin{example}[: the AI at 08:30 on 25 August]
The network says $p=0.75$: above 0.25, so ``raining''. Its correction is $c = 4.00$\,dB (at its upper limit; the
classical wet-antenna model said 4.10\,dB). So
\[
A = 9.48 - 4.00 = 5.48\ \text{dB},\qquad
\frac{A}{aL} = \frac{5.48}{0.370\times 1.066} = 13.9,
\]
\[
R = 13.9^{1/0.860} \approx 21.3\ \text{mm/h}.
\]
The classical method gave 20.7\,mm/h for the same minute. Averaged over the 5 minutes around 08:30, the AI says
16.4\,mm/h, the classical method 17.3 and the gauge 14.4.
\end{example}

\begin{figure}[h]
\centering
\begin{tikzpicture}
\begin{groupplot}[group style={group size=1 by 2,vertical sep=0.9cm,xlabels at=edge bottom},
  width=0.92\textwidth,height=4.8cm,xmin=4,xmax=24,xtick={4,6,...,24},grid=major,grid style={inkc!12},
  tick label style={font=\small},label style={font=\small},legend style={font=\scriptsize,draw=none,fill=white,fill opacity=0.85,text opacity=1},legend cell align=left]
\nextgroupplot[ylabel={rain (mm/h)},ymin=0,ymax=24,legend pos=north east]
  \addplot[thick,color=gaugec,forget plot] table[x=hour,y=gauge] {storm.dat};
  \addplot[thick,color=rainc!55,forget plot] table[x=hour,y=R] {storm.dat};
  \addplot[thick,color=inkc,forget plot] table[x=hour,y=ai] {storm_ai.dat};
  \addlegendimage{thick,color=gaugec} \addlegendentry{rain gauge}
  \addlegendimage{thick,color=rainc!55} \addlegendentry{classical method (Part 1)}
  \addlegendimage{thick,color=inkc} \addlegendentry{AI}
\nextgroupplot[ylabel={dB},xlabel={hour of the day (UTC), 25 August 2015},ymin=-1.5,ymax=10,legend pos=north east]
  \addplot[very thin,color=rainc,forget plot] table[x=hour,y=excess] {storm_ai.dat};
  \addplot[thick,color=gaugec,forget plot] table[x=hour,y=corr] {storm_ai.dat};
  \addlegendimage{thick,color=rainc} \addlegendentry{extra loss (over the 3-day dry level)}
  \addlegendimage{thick,color=gaugec} \addlegendentry{AI's correction $c$}
\end{groupplot}
\end{tikzpicture}
\caption{The same storm as Figure~\ref{fig:storm}, now with the AI. Top: rain from the AI (5-minute averages), the
classical method and the gauge 4\,m away. Bottom: what the AI saw (extra loss) and how much of it the AI decided was
not rain in the air. In dry spells the correction stays near 0\,dB; there the rain/no-rain answer keeps the rain at
zero. In the storm the correction climbs to its 4\,dB limit (wet antennas), and only the loss above that is turned into
rain. The AI under-reads the short burst at 17:05 (about 4\,mm/h against the gauge's 12).}
\label{fig:storm_ai}
\end{figure}

\subsection{How it was trained}
\begin{itemize}
\item \textbf{Examples:} the 683 Gothenburg sub-links that passed quality control in Phase~1, over the training
weeks. The ``right answer'' is the weather radar's rain along each link path, available every 5 minutes.
\item \textbf{One round:} 64 pieces of signal, each 2{,}560 minutes long (about 43 hours), picked at random links and
random times. The first 1{,}024 minutes of each piece are only there to fill the network's 17-hour view; they are not
scored.
\item \textbf{Penalties:} the rain/no-rain penalty on every radar minute, and the squared-mistake penalty for the
amount \emph{only on minutes where the radar says it rained}. (Why only those? See ``What went wrong first'' below.)
\item \textbf{Practice with coarse data:} in half the rounds, the signal was first rounded to whole decibels, the way
much operator equipment stores it (Italy's does). This is called \textbf{augmentation}: showing the same examples in
a slightly different form so the network learns not to depend on details that change between networks.
\item \textbf{How long:} 4{,}000 rounds, about one minute on one graphics card (GPU, a chip built for doing many
multiplications at once). Every 250 rounds we checked the validation fortnight and kept the best version.
\item \textbf{Five tries, keep one:} training starts from random dials, and two runs never end up identical (more on this
in Section~\ref{sec:se_results}). So we trained five times and kept the run with the smallest penalty on the validation
fortnight. That choice uses no test data. The kept network has 111{,}890 dials.
\end{itemize}

\subsection{What went wrong first}
Our first designs failed, and how they failed is instructive.
\begin{enumerate}
\item \textbf{Letting the network say the mm/h directly.} It caught rain well, but reported only about half the rain.
The training weeks (June, early July) had light rain; August had heavy storms. When the network meets a signal drop
bigger than anything it practised on, it plays safe and answers something it has seen before. A network that only
learned from examples has no reason to believe the rain can be heavier than in its examples.
\item \textbf{Physics conversion, but an unlimited correction learned on all minutes.} Now it reported a quarter of the
rain. On the many dry minutes it learned that subtracting a lot is always safe, and it carried the habit into storms.
\item \textbf{A dry level that fails after a rainy day.} The 24-hour median of Section~\ref{sec:drylevel}. This did not
show up in Sweden at all; it showed up in Italy and the Netherlands.
\end{enumerate}
The fixes are the design described above: a correction limited to $-1\ldots4$\,dB, the amount learned only when it
rains, and the 3-day low dry level. The physics conversion extrapolates to heavy rain by itself, because the power law
does not care whether the network practised on heavy storms.

\begin{careful}
we made the first two fixes after looking at the Swedish August scores, and the third after looking at the Italian and
Dutch results. So none of the three tests is a perfectly untouched exam any more: in each case we chose part of the
design after seeing how it did there. What we did \emph{not} do is tune numbers on them: the 3-day window and the 20\%
were chosen from reasoning, and nothing was trained on Italian or Dutch data. A truly fresh test needs data we have not
looked at yet, such as another year or another country.
\end{careful}

%==========================================================
\section{Before the results: exactly how we score}
\label{sec:scoring_rules}
%==========================================================

Section~\ref{sec:scores} introduced the scores. Before using them on three countries, here are the exact rules, because
small choices in how scores are counted can make a method look better or worse than it is. Both methods, classical and AI,
are always scored with the same rules, on the same links and the same hours.

\subsection{Two separate exams}
Every test has two parts, and they answer different questions:
\begin{itemize}
\item \textbf{Rain or no rain?} Only the yes/no answer is checked, not the amount. Scores: MCC, share of rain caught,
share of false alarms.
\item \textbf{How much rain?} The method's millimetres per hour are compared with the measured millimetres per hour. Scores:
correlation $r$, typical error (RMSE) and bias.
\end{itemize}
A method can pass one and fail the other. A detector that always says ``raining'' during storms but gets the amount
wrong by half would score well on the first and badly on the second.

\subsection{Rule 1: what counts as ``rain''}
A time step counts as \textbf{rainy} if the reference measured \textbf{at least 0.1\,mm/h} there. Less than that is
too little to matter and is often just noise in the reference itself.

Our methods give an answer every minute, but the references come every 5, 15 or 60 minutes. So for the yes/no exam we
look at the minutes inside each reference period and call the period rainy for the method if \textbf{more than half}
of those minutes were rainy.

\subsection{Rule 2: dry moments are left out of the amount scores}
For the amounts, a time step is only scored if \textbf{the reference or the method} shows at least 0.1\,mm/h. Moments
where both agree it was dry are skipped. Why? Most hours are dry. Every dry hour where both say 0 is a ``perfect''
answer that has nothing to do with measuring rain, and thousands of them would make the typical error look tiny.

\begin{example}[: why dry hours are skipped]
Five hours, reference against method (mm/h):
\begin{center}\small
\begin{tabular}{lccccc}
\toprule
Hour & 1 & 2 & 3 & 4 & 5\\
\midrule
Reference & 0 & 0 & 2.0 & 8.0 & 0.05\\
Method & 0 & 0.3 & 1.5 & 6.0 & 0\\
\midrule
Scored? & no (both dry) & yes (method $\geq$ 0.1) & yes & yes & no (both below 0.1)\\
\bottomrule
\end{tabular}
\end{center}
The three scored mistakes are $+0.3$, $-0.5$ and $-2.0$\,mm/h. Then
\[
\text{RMSE} = \sqrt{\frac{0.3^2 + 0.5^2 + 2.0^2}{3}} = \sqrt{\frac{0.09+0.25+4}{3}} = 1.20\ \text{mm/h},
\]
\[
\text{bias} = \frac{(0.3+1.5+6.0)-(0+2.0+8.0)}{0+2.0+8.0} = -22\%.
\]
Had we scored all five hours, the RMSE would have dropped to 0.93\,mm/h, only because of the two dry hours, though the
method measured the rain no better.
\end{example}

Notice that hour 2 \emph{is} scored: when the method claims rain the reference did not see, that false rain counts as a
mistake. A method cannot hide its false alarms by skipping them.

\subsection{Rule 3: which reference, and how often}
\begin{center}\small
\begin{tabular}{p{2.3cm}p{7.4cm}p{3.9cm}}
\toprule
Where & Reference & Compared every\\
\midrule
Sweden & weather radar, averaged along the link & 5 minutes and 1 hour\\
       & rain gauges within 1\,km of the link & 1 hour\\
Italy & gauge-adjusted radar, averaged along the link & 15 minutes and 1 hour\\
      & rain gauges within 1\,km of the link & 1 hour\\
Netherlands & gauge-adjusted radar, averaged along the link & 1 hour\\
\bottomrule
\end{tabular}
\end{center}
``Averaged along the link'' means: take every radar square the beam crosses and average them, giving more weight to the
squares where the beam travels a longer distance. That is the same thing the link itself measures.

\subsection{No reference is the true rain}
It is tempting to call the gauge ``the truth''. It is not, for this job:
\begin{itemize}
\item \textbf{A gauge measures one spot}, the size of a bucket. The link measures the average along a line that can be
several kilometres long. A storm cell 2\,km wide can soak the gauge and only half the link, or the other way round. So
even a perfect link and a perfect gauge disagree. That is why we only pair a link with gauges within 1\,km of it, and why
we compare hours rather than minutes (over an hour, a passing storm crosses both).
\item \textbf{Radar sees the whole line} but measures echoes high in the air and converts them to rain with a formula of
its own. Plain radar (Sweden) can be off by tens of percent; gauge-adjusted radar (Italy, the Netherlands) has been
corrected with gauges and is better, but still averages over 1\,km squares.
\end{itemize}
So a correlation of 1 is impossible even for a perfect method. As a yardstick: in Sweden, the radar itself agrees with the
gauges at $r = 0.80$ (hourly). A link method scoring around 0.8 against gauges is as close as the radar gets.

\begin{oneline}
We score two things separately (rain or not, and how much), count rain from 0.1\,mm/h, leave out moments where everyone
agrees it was dry, and compare with references that are good but not perfect.
\end{oneline}

%==========================================================
\section{Results in Sweden}
\label{sec:se_results}
%==========================================================

Same links, same month (August 2015), same scores as Section~\ref{sec:scores}.

\begin{center}\small
\begin{tabular}{lcccc}
\toprule
& \multicolumn{2}{c}{\textbf{Rain or no rain}} & \multicolumn{2}{c}{\textbf{How much (hourly, vs gauges)}}\\
\cmidrule(lr){2-3}\cmidrule(lr){4-5}
Method & MCC & caught / false alarms & correlation $r$ & error (RMSE) / bias\\
\midrule
Classical (Part 1) & 0.578 & 71\% / 4.6\% & 0.81 & 1.04\,mm/h / $-6\%$ \\
AI & \textbf{0.615} & 68\% / \textbf{3.2\%} & 0.79 & 1.04\,mm/h / $-20\%$ \\
\midrule
\multicolumn{5}{l}{\emph{Same August data, but rounded to whole decibels first:}}\\
Classical & 0.336 & 82\% / 23\% & 0.85 & 0.82\,mm/h / $+26\%$ \\
AI & \textbf{0.568} & 65\% / \textbf{3.9\%} & 0.80 & 0.96\,mm/h / \textbf{$-15\%$} \\
\bottomrule
\end{tabular}
\end{center}

Reading the table:
\begin{itemize}
\item \textbf{Rain or no rain:} the AI is better. It raises fewer false alarms (3.2\% of dry periods instead of
4.6\%) at the cost of missing slightly more light rain. And it does this without looking into the future.
\item \textbf{Amounts:} the classical method is better. The AI follows the ups and downs about as well (correlation 0.79
against 0.81, same typical error), but it measures a fifth too little rain, against 6\% too little for the classical
method. Against the radar (not in the table) it is also slightly behind: hourly $r$ 0.71 against 0.76.
\item \textbf{Coarse data:} this is the clearest win. Rounded to whole decibels, the classical detector panics (23\%
false alarms, because 1-dB steps look like jumpiness), while the AI barely changes.
\end{itemize}

\subsection{Is the AI just lucky?}
Training starts from random dials, so two training runs never give exactly the same network. Here are all five training
runs of the same design, with different random starts (``seeds''):
\begin{center}\small
\begin{tabular}{lcc}
\toprule
& average of 5 runs & worst -- best \\
\midrule
MCC & 0.619 & 0.612 -- 0.637 \\
hourly correlation with gauges & 0.79 & 0.77 -- 0.84 \\
hourly error vs gauges (RMSE) & 1.05\,mm/h & 0.90 -- 1.15 \\
bias vs gauges & $-22\%$ & $-32\%$ -- $-6\%$ \\
\bottomrule
\end{tabular}
\end{center}
So: the better rain detection is solid (every run beats 0.578). The amounts are not: one run under-measures by only
6\%, another by 32\%, with nothing but the random start to tell them apart. The run we kept (chosen on validation, not on
these numbers) sits in the middle, at $-20\%$. Picking the run that happens to score best in August would have made the
AI look much better than it is. Section~\ref{sec:why_runs} finds why the runs differ: not in \emph{when} they say it rains,
but in how strongly they turn a decibel into millimetres, which each run learns from a handful of heavy Swedish minutes.
Training on more varied rain brings them together (Section~\ref{sec:cv}).

\subsection{Which design choices mattered?}
Removing one ingredient at a time (an \textbf{ablation}) shows what each does:
\begin{center}\footnotesize
\begin{tabular}{p{6.3cm}cccc}
\toprule
Version (one training run each) & MCC & $r$ vs gauges & bias vs gauges & MCC, 1-dB data\\
\midrule
Full design & 0.615 & 0.79 & $-20\%$ & 0.568\\
Network says mm/h directly (no physics conversion) & 0.618 & 0.75 & $-19\%$ & 0.550\\
No practice on 1-dB rounded data & 0.634 & 0.79 & $-27\%$ & 0.481\\
\bottomrule
\end{tabular}
\end{center}
The practice on rounded data clearly matters for coarse data (0.481 without it, against 0.568). For the physics
conversion the picture is less clear than we first thought: with the final design, a network that outputs mm/h
directly does about as well \emph{in Sweden}. Differences this small are within the spread between training runs shown
above. The physics conversion was what rescued our first designs from heavy under-measurement; whether it still helps
abroad, we have not tested.

\begin{oneline}
In Sweden the AI detects rain better than the classical method, works live, and copes far better with coarse data, but
it measures about a fifth too little rain, and how much too little changes from one training run to the next.
\end{oneline}

%==========================================================
\section{A harder exam: Italy}
\label{sec:italy}
%==========================================================

A method that only works in the city where it was trained is of little use for Algeria. So we took the network trained
in Gothenburg and, \textbf{without any change}, ran it on a completely different network: the Lepida network in
Emilia-Romagna, northern Italy (the OpenRainER dataset). Using a model on a place it never learned from is called a
\textbf{zero-shot} test. The classical method was also used with its Swedish settings.

Almost everything is different there:
\begin{center}\small
\begin{tabular}{p{3.4cm}p{4.6cm}p{5.6cm}}
\toprule
& Sweden (training) & Italy (test)\\
\midrule
Frequencies & 7.5 -- 39\,GHz & only 24.5 -- 25.6\,GHz\\
Link lengths & half under 2.3\,km, up to 15\,km & up to 25\,km (after quality control: up to 9\,km)\\
Landscape & flat coastal city & plains, hills and Apennine mountains\\
Rounding of readings & 0.1\,dB & 1\,dB\\
Period & summer 2015 & May -- September 2021 and 2022\\
\bottomrule
\end{tabular}
\end{center}

The references in Italy are good: a \textbf{gauge-adjusted radar} (radar maps corrected with hundreds of rain gauges,
every 15 minutes, 1\,km squares) and 22--23 gauges standing within 1\,km of a link. One detail had to be checked: the files
do not say whether a 15-minute value is labelled at the start or the end of its 15 minutes. We tried both on the classical
estimates; ``end'' fits clearly better (correlation 0.25 against 0.18), and we used it for both methods.

\begin{careful}
we applied the same quality-control rule as in Sweden (drop a sub-link if it has less than half its data or is flagged
``wet'' more than 30\% of the time by the classical detector). With 1-dB readings the classical detector flags many
links as almost always wet, so only 81 of 302 sub-links (2021) and 72 of 302 (2022) survived. Both methods are scored
on exactly the same surviving links, but the rule judges links with one of the two methods being compared.
Section~\ref{sec:neighbours} tries the fairer alternative: keep every link that has data, and let a neighbour check
catch the bad ones.
\end{careful}

\begin{center}\footnotesize
\begin{tabular}{llcccc}
\toprule
Season & Method & MCC & caught / false alarms & vs radar (hourly): $r$ / bias & vs gauges (hourly): $r$ / bias\\
\midrule
2021 & classical & 0.25 & 75\% / 16.7\% & 0.22 / $+64\%$ & 0.45 / $+39\%$\\
     & AI        & \textbf{0.51} & 57\% / \textbf{1.8\%} & \textbf{0.70} / $\mathbf{-5\%}$ & \textbf{0.74} / $\mathbf{-3\%}$\\
2022 & classical & 0.27 & 70\% / 16.3\% & 0.12 / $+132\%$ & 0.31 / $+64\%$\\
     & AI        & \textbf{0.50} & 51\% / \textbf{1.9\%} & \textbf{0.68} / $\mathbf{-7\%}$ & \textbf{0.81} / $\mathbf{-10\%}$\\
\bottomrule
\end{tabular}
\end{center}

The AI wins clearly on both exams. The classical detector raises false alarms on one dry period in six, the AI on one in
fifty. For amounts, the AI's hourly correlation with the gauges (0.74 and 0.81) is as good as in Sweden, and its total is
within 10\% of the truth, while the classical method's correlations fall to 0.12--0.45 and it measures far too much.
Why does the classical method collapse?

\subsection{A few broken hours spoil the average}
The correlation in the table is computed over all links and hours together (``pooled''). Correlation is very sensitive to
a few wild values: one hour where a link claims 400\,mm/h while the radar says 2 can drag it down more than a thousand good
hours lift it up. So we also computed it the other way: for \emph{each link on its own}, how well do its hours follow the
radar? And we counted the impossible hours (more than 100\,mm/h for a whole hour, which essentially never happens there).

\begin{center}\small
\begin{tabular}{llcc}
\toprule
Season & Method & typical link's $r$ (median) & impossible hours ($>$100\,mm/h)\\
\midrule
2021 & classical & 0.84 & 35\\
     & AI        & 0.82 & 0\\
2022 & classical & 0.79 & 124\\
     & AI        & 0.79 & 0\\
\midrule
\multicolumn{2}{l}{Sweden, for comparison (classical / AI)} & 0.88 / 0.87 & \\
\bottomrule
\end{tabular}
\end{center}

This explains it. \textbf{The typical Italian link follows the radar almost as well as a Swedish one} (about 0.8
against 0.88), with both methods. What ruins the classical pooled numbers is a small number of hours where a link
misbehaves, for example a sudden drop of the signal caused by equipment rather than rain, read as a cloudburst. The
AI produced none at all.

An earlier version of our network, which had only the 24-hour median as its dry level, still produced 8 and 19 such
hours, and its pooled correlation was only 0.33. Adding the second dry level removed them. We think we know why, but
have not proven it: a signal that jumps to a new level and stays there (equipment) moves the 24-hour median and the
3-day low value very differently from rain, which comes and goes. With both in front of it, the network can tell the
two apart.

\begin{figure}[h]
\centering
\begin{tikzpicture}
\begin{axis}[width=0.95\textwidth,height=5.2cm,xmin=11,xmax=14.5,ymin=0,ymax=36,grid=major,grid style={inkc!12},
  xtick={11,11.5,...,14.5},xticklabels={11:00,11:30,12:00,12:30,13:00,13:30,14:00,14:30},
  xlabel={time of day (UTC), 8 June 2021},ylabel={rain (mm/h, 15-min average)},
  tick label style={font=\small},label style={font=\small},
  legend style={font=\scriptsize,draw=none,fill=white,fill opacity=0.85,text opacity=1},legend cell align=left,legend pos=north east]
  \addplot[thick,color=gaugec,mark=*,mark size=1.3pt] table[x=hour,y=gauge] {italy_storm.dat};
  \addlegendentry{rain gauge, 84\,m from the link}
  \addplot[thick,color=radarc,mark=square*,mark size=1.3pt] table[x=hour,y=radar] {italy_storm.dat};
  \addlegendentry{gauge-adjusted radar along the link}
  \addplot[thick,color=rainc!55,mark=triangle*,mark size=1.6pt] table[x=hour,y=classical] {italy_storm.dat};
  \addlegendentry{classical method}
  \addplot[thick,color=inkc,mark=diamond*,mark size=1.6pt] table[x=hour,y=ai] {italy_storm.dat};
  \addlegendentry{AI (trained in Sweden only)}
\end{axis}
\end{tikzpicture}
\caption{A short cloudburst in Italy (sub-link 154, 2.1\,km, 24.6\,GHz). This storm was picked by a fixed rule (the
closest link--gauge pair that saw more than 20\,mm/h), not because it looks good. Both link methods are shown after the
neighbour check of Section~\ref{sec:neighbours}. Storm totals: gauge 11.0\,mm, AI 6.8\,mm, classical 6.2\,mm, radar
4.4\,mm. A storm lasting half an hour over a few square kilometres is exactly the kind
radar smooths away and flash floods come from.}
\label{fig:italy}
\end{figure}

%==========================================================
\section{A different kind of data: the Netherlands}
\label{sec:nl}
%==========================================================

\subsection{Minimum and maximum instead of averages}
Gothenburg's links report the signal every 10 seconds, and Italy's every minute. Many operators keep far less. The Dutch
operator, like many, stored only \textbf{two numbers per 15 minutes}: the weakest and the strongest received level in
that quarter of an hour. This is called \textbf{min/max data}.

\begin{figure}[h]
\centering
\begin{tikzpicture}[x=0.6cm,y=0.5cm]
  \draw[inkc!40] (0,0) rectangle (15,6);
  \draw[rainc,thick] plot[smooth,tension=0.5] coordinates {(0,4.6) (1,4.7) (2,4.4) (3,4.5) (4,3.9) (5,3.2) (6,2.1) (7,1.4) (8,1.9) (9,1.2) (10,2.4) (11,3.1) (12,3.8) (13,4.2) (14,4.5) (15,4.6)};
  \draw[gaugec,dashed] (0,4.75) -- (15,4.75) node[right,font=\scriptsize] {strongest = max};
  \draw[gaugec,dashed] (0,1.15) -- (15,1.15) node[right,font=\scriptsize] {weakest = min};
  \node[font=\scriptsize,anchor=north] at (7.5,0) {15 minutes};
  \node[font=\scriptsize,rotate=90,anchor=south] at (0,3) {received level};
\end{tikzpicture}
\caption{What min/max data keeps. The real signal (blue) wanders during 15 minutes; only its weakest and strongest
values are stored. Everything in between is thrown away.}
\label{fig:minmax}
\end{figure}

Our network was built for minute-by-minute data, so it cannot read this directly. But we can turn the Swedish data into
the Dutch format: for every 15 minutes, take the weakest and strongest of the ninety 10-second readings (and, in half of
the practice rounds, round them to whole decibels first, like the Dutch Nokia equipment). We trained a min/max version of
the same network on that, with the same weeks and the same rules as before. Its view reaches back 63 hours.

For comparison we built a classical min/max method in the style of \textbf{RAINLINK}, the algorithm developed at the Dutch
meteorological institute (KNMI) and Wageningen University for this kind of data: call it rain if the weakest level falls more than 1.75\,dB below the
24-hour reference, subtract a fixed 2.3\,dB for wet antennas, and mix the rain from the weakest and strongest readings
as $0.33\times$ (from weakest) $+ 0.67\times$ (from strongest). The 1.75\,dB and the 2.3\,dB were chosen on Swedish June--July
data. (Real RAINLINK also compares each link with its neighbours to throw out links that drop alone; we test our own
version of that, for both methods, in Section~\ref{sec:neighbours}.)

\subsection{The test}
We then ran both, unchanged, on the Netherlands from 1 June to 22 August 2014: 1{,}081 sub-links (after the same quality
rule, plus dropping links below 10\,GHz, which rain barely affects), 15--58\,GHz, 0.1 to 16\,km long. The reference is
KNMI's gauge-adjusted radar, one value per hour per 1\,km square, averaged along each link. Polarisation is not in the
Dutch files, so we assumed vertical.

\begin{center}\footnotesize
\begin{tabular}{lccccc}
\toprule
Method & MCC (hourly) & $r$ (pooled) & RMSE & bias & typical link's $r$ / impossible hours\\
\midrule
Classical min/max & \textbf{0.63} & 0.16 & 8.7\,mm/h & $+54\%$ & 0.79 / 219\\
AI min/max        & 0.55 & \textbf{0.50} & \textbf{2.5\,mm/h} & \textbf{$+16\%$} & 0.77 / \textbf{26}\\
\bottomrule
\end{tabular}
\end{center}

Here the classical method detects rain better (hourly MCC 0.63 against 0.55). For amounts, the same story as Italy: the
typical link works with both methods (0.77--0.79), the classical method produces many impossible hours that destroy its
pooled score, and the AI produces eight times fewer, though not none. Split by equipment maker, the AI does better on
Nokia links (1-dB readings, bias $+10\%$) than on NEC links (0.1-dB readings, $+41\%$).

\begin{figure}[h]
\centering
\begin{tikzpicture}
\begin{axis}[width=0.98\textwidth,height=5.4cm,xmin=-1,xmax=83,ymin=0,ymax=26,
  xlabel={day (1 June -- 22 August 2014)},ylabel={rain (mm/day)},
  xtick={0,30,61},xticklabels={1 June,1 July,1 August},grid=major,grid style={inkc!12},
  tick label style={font=\small},label style={font=\small},
  legend style={font=\scriptsize,draw=none,fill=white,fill opacity=0.85,text opacity=1},legend cell align=left,legend pos=north west]
  \addplot[ycomb,line width=2.2pt,color=radarc!55,forget plot] table[x=day,y=radar] {nl_daily.dat};
  \addlegendimage{line width=2.2pt,color=radarc!55}
  \addlegendentry{KNMI radar}
  \addplot[only marks,mark=triangle*,mark size=1.8pt,color=rainc!70] table[x=day,y=classical] {nl_daily.dat};
  \addlegendentry{classical min/max}
  \addplot[only marks,mark=*,mark size=1.4pt,color=inkc] table[x=day,y=ai] {nl_daily.dat};
  \addlegendentry{AI min/max}
\end{axis}
\end{tikzpicture}
\caption{Daily rain averaged over every Dutch link with enough data (1{,}908 sub-links), after the neighbour check of
Section~\ref{sec:neighbours}. Day to day both methods follow the radar well (correlation 0.90 classical, 0.94 AI). Over the
whole summer the classical method adds up to 223\% of the radar's rain and the AI to 132\%: with every link included,
both measure too much, the classical method far more.}
\label{fig:nl}
\end{figure}

\begin{oneline}
Trained only in Sweden, the AI ran in Italy and the Netherlands with no retraining; the typical link worked in both
countries, the AI made no absurd readings in Italy and far fewer than the classical method in the Netherlands, where a guard
against misbehaving links is still needed.
\end{oneline}

%==========================================================
\section{Catching misbehaving links: the neighbour check}
\label{sec:neighbours}
%==========================================================

So far, a small number of links that misbehave, for example a signal that suddenly drops because of the equipment, have
spoiled the averages in Italy and the Netherlands. And in Italy we kept only the links the classical detector liked. Both
problems have the same natural fix, used by the Dutch RAINLINK method: \textbf{ask the neighbours}.

\subsection{The idea}
Rain falls from clouds that are kilometres wide. If it really rains on one link, other links close by almost always feel it
too, within minutes. A link that claims a storm while every link around it stays dry is far more likely broken than wet.

Our rule, applied in exactly the same way to the classical method and to the AI:
\begin{enumerate}
\item \textbf{Neighbours} of a link are the other links whose middle point is within \textbf{10\,km} of its own. (The two
directions of one physical link do not count as each other's neighbours: they share the same equipment and the same air,
so they would fail together.)
\item A link's rain counts only if at least \textbf{one neighbour also reports rain} (0.1\,mm/h or more) within
\textbf{15 minutes before or after}. Otherwise it is set to zero.
\item If \emph{no} neighbour has any data at that moment, the link cannot be checked and is left alone.
\end{enumerate}
We fixed 10\,km and 15 minutes from the size and speed of storms (a storm moving at 40\,km/h crosses 10\,km in 15 minutes),
before looking at any results. The price is a delay: a warning may need up to 15 minutes to be confirmed by a neighbour.

\begin{example}[: three links, one storm, one fault]
Links A and B are 4\,km apart; link C is 6\,km from both. At 14:00 link A reports 12\,mm/h. Link B reports rain from 14:10.
Link A's rain counts: B confirms it within 15 minutes. At 20:00 link C reports 40\,mm/h for an hour while A and B, both
working, report nothing from 19:45 to 21:15. C's rain is set to zero: no neighbour confirms it.
\end{example}

\subsection{Results on the same links as before}
First on exactly the links of the previous sections, to see what the check alone does. ``Impossible hours'' are hours with
more than 100\,mm/h, as in Section~\ref{sec:italy}.
\begin{center}\footnotesize
\begin{tabular}{llccc}
\toprule
Test (links) & Method & MCC & impossible hours & bias vs gauges / radar\\
\midrule
Sweden (683) & classical & 0.578 $\to$ 0.581 & 0 $\to$ 0 & $-6\%$ $\to$ $-6\%$\\
             & AI        & 0.617 $\to$ 0.622 & 0 $\to$ 0 & $-19\%$ $\to$ $-19\%$\\
Italy 2021 (81) & classical & 0.251 $\to$ 0.254 & 35 $\to$ 26 & $+39\%$ $\to$ $+28\%$\\
                & AI        & 0.509 $\to$ 0.518 & 0 $\to$ 0 & $-3\%$ $\to$ $-7\%$\\
Italy 2022 (72) & classical & 0.269 $\to$ 0.271 & 124 $\to$ 77 & $+64\%$ $\to$ $+36\%$\\
                & AI        & 0.502 $\to$ 0.506 & 0 $\to$ 0 & $-10\%$ $\to$ $-14\%$\\
Netherlands (1{,}081) & classical & 0.626 $\to$ 0.632 & 219 $\to$ 163 & $+54\%$ $\to$ $+43\%$\\
                      & AI        & 0.547 $\to$ 0.562 & 26 $\to$ 16 & $+16\%$ $\to$ $+12\%$\\
\bottomrule
\end{tabular}
\end{center}
(Bias is against gauges for Sweden and Italy, against radar for the Netherlands. The AI's Swedish MCC is 0.617 here rather
than 0.615 because it is scored only where the link has data, like the classical method.)

The check helps both methods everywhere, a little: rain detection never gets worse, and impossible hours drop by a quarter
to a half. It does not remove them all. A broken classical link is often ``confirmed'' by a neighbour that is itself
raising a false alarm, which happens easily when the classical detector raises false alarms on one dry period in six. In
Sweden, where links are dense (a typical link has 192 neighbours) and well behaved, the check removes under 2\% of rainy
minutes and changes almost nothing. In Italy, where a typical link has only 6--7 neighbours, it removes about a fifth of the
rainy minutes; the AI's amounts drop by only about 5\%, so most of what it removed was little or no rain, but some real rain
went too.

\subsection{Keeping every link}
Now the fairer filter: keep \textbf{every link that has data at least half of the time}, and let the neighbour check do the
cleaning. In Italy that means 281 and 267 sub-links instead of 81 and 72; in the Netherlands 1{,}908 instead of 1{,}081
(still dropping links below 10\,GHz).
\begin{center}\footnotesize\setlength{\tabcolsep}{4pt}
\begin{tabular}{llcccc}
\toprule
Test (links) & Method, with neighbour check & MCC & impossible hours & $r$ vs gauges & bias vs gauges\\
\midrule
Italy 2021 (281) & classical & 0.158 & 43 & 0.45 & $+76\%$\\
                 & AI        & \textbf{0.449} & \textbf{0} & \textbf{0.72} & $\mathbf{+3\%}$\\
Italy 2022 (267) & classical & 0.168 & 106 & 0.38 & $+49\%$\\
                 & AI        & \textbf{0.443} & \textbf{5} & \textbf{0.69} & $-16\%$\\
\midrule
& & & & $r$ vs radar & bias vs radar\\
\midrule
Netherlands (1{,}908) & classical & 0.463 & 199 & 0.19 & $+128\%$\\
                      & AI        & \textbf{0.484} & \textbf{19} & \textbf{0.52} & $\mathbf{+34\%}$\\
\bottomrule
\end{tabular}
\end{center}

Two things stand out:
\begin{itemize}
\item \textbf{The extra links are worse links.} With every link included, both methods score lower than on the hand-picked
set (Italian rain detection 0.45 instead of 0.52 for the AI; Dutch 0.48 instead of 0.56). The old filter was unfair, but it
was throwing out links that deserved it. The neighbour check alone does not replace it.
\item \textbf{The AI copes much better with bad links.} On every link, it still detects rain about three times better than
the classical method in Italy, has almost no impossible hours, and stays within 3--16\% of the Italian gauges, while the
classical method measures half to three-quarters too much. In the Netherlands the AI now also beats the classical method at
detecting rain (0.48 against 0.46).
\end{itemize}
For a real operator network, where nobody hand-picks the links, this second table is the honest picture. What it calls for
is a fair way to judge each link from its own behaviour, instead of throwing out links with the classical detector.

\subsection{A reliability score for each link}
\label{sec:reliability}
The score looks only at the raw signal, never at the rain worked out by either method, so it cannot favour one of them.
Every 15 minutes we compare the link's total loss with its own dry level (the 3-day low value of
Section~\ref{sec:drylevel}). If it is at least \textbf{1.5\,dB} above, we call that 15 minutes a \textbf{signal event}:
something is weakening the beam. Then two questions about each link:
\begin{enumerate}
\item \textbf{Does it rain alone?} For each of its events, did at least one of its \textbf{5 nearest neighbours} (within
10\,km) also have an event within 15 minutes? If more than half of its events had no such company (and it had at least 8
events, enough to judge), the link is dropped.
\item \textbf{Is it always ``raining''?} If the link is in an event more than \textbf{20\%} of the time, it is dropped. In
these climates it rains far less than a fifth of the time; a beam that is weakened that often is drifting or faulty.
\end{enumerate}
Links with at least half their data that pass both questions are kept.

\begin{example}[: one link judged]
Over a summer, a link has 40 signal events. For 30 of them at least one of its five nearest neighbours also had an event
within 15 minutes; 10 had no company. Unconfirmed share: $10/40 = 25\%$, below one half, so the first question passes. Its
40 events are 40 of its roughly 14{,}000 quarter-hours of data, 0.3\% of the time, far below 20\%. The link is kept. A link
whose signal drops every other evening while its neighbours stay calm would have most of its events unconfirmed and be
dropped.
\end{example}

\textbf{How we chose the numbers.} We set them on Swedish June--July data only, where we could check them against the radar
without touching any test month. Our first version asked ``any neighbour within 10\,km''. In Gothenburg a typical link has
192 neighbours that close, and by pure chance one of them is almost always having an event, so almost nothing was ever
judged lonely. Asking only the 5 nearest, and adding the 20\% question, fixed that. On the Swedish June--July data, the
final score drops 69 of 728 sub-links. It agrees with the old rule on 40 of the 45 that rule dropped, and drops 29 more.
Are those 29 really bad? Their raw signal follows the radar with a typical correlation of 0.33, as poorly as the links
both rules drop (0.32), against 0.58 for the links kept. So yes.

\textbf{Results.} The same three link sets, both methods with the neighbour check:
\begin{center}\footnotesize\setlength{\tabcolsep}{4pt}
\begin{tabular}{lllccc}
\toprule
Test & Links & Method & MCC & $r$ / bias$^{*}$ & impossible hours\\
\midrule
Italy 2021 & old rule (81) & AI & 0.518 & 0.75 / $-7\%$ & 0\\
           & all with data (281) & AI & 0.449 & 0.72 / $+3\%$ & 0\\
           & \textbf{reliability score (196)} & AI & \textbf{0.521} & \textbf{0.78} / $-6\%$ & 0\\
           & reliability score (196) & classical & 0.186 & 0.47 / $+47\%$ & 37\\
\midrule
Italy 2022 & old rule (72) & AI & 0.506 & 0.84 / $-14\%$ & 0\\
           & all with data (267) & AI & 0.443 & 0.69 / $-16\%$ & 5\\
           & \textbf{reliability score (194)} & AI & \textbf{0.508} & 0.72 / $-20\%$ & 2\\
           & reliability score (194) & classical & 0.193 & 0.44 / $+23\%$ & 98\\
\midrule
Netherlands & old rule (1{,}081) & AI & 0.562 & 0.53 / $+12\%$ & 16\\
            & all with data (1{,}908) & AI & 0.484 & 0.52 / $+34\%$ & 19\\
            & \textbf{reliability score (1{,}653)} & AI & 0.508 & \textbf{0.55} / $+20\%$ & 18\\
            & reliability score (1{,}653) & classical & 0.517 & 0.22 / $+52\%$ & 180\\
\bottomrule
\end{tabular}

\smallskip
$^{*}$ against gauges in Italy, against the radar in the Netherlands.
\end{center}

\begin{itemize}
\item \textbf{Many more links, same quality, in Italy.} The score keeps 196 and 194 Italian sub-links, more than twice as
many as the old rule, and the AI detects rain as well as on the hand-picked set (0.52 and 0.51). In 2021 it even agrees
better with the gauges (0.78). In 2022 the gauge agreement stays below the old set's (0.72 against 0.84): some of the
extra links are still weaker.
\item \textbf{A middle way in the Netherlands.} 1{,}653 sub-links, in between ``everything'' and the old set, and in between
on rain detection too; the hourly correlation with the radar is the best of the three.
\item \textbf{It does not rescue the classical method.} On the reliable links the classical method still raises too many
false alarms in Italy (MCC 0.19) and keeps many impossible hours. Its problem there is its detector on 1-dB data, not
bad links.
\end{itemize}

\begin{oneline}
Checking each link's rain against its neighbours helps both methods; judging each link by whether its signal drops alone or
all the time lets us keep two to three times more Italian links than hand-picking, with the AI scoring as well as before.
\end{oneline}

%==========================================================
\section{Teaching the AI with more varied weather}
\label{sec:varied}
%==========================================================

Two questions were left open above. Why does the AI still measure too little rain after a rainy day in Italy
(Section~\ref{sec:drylevel})? And why do training runs that agree on \emph{whether} it rains disagree so much on
\emph{how much} (Section~\ref{sec:se_results})? We first looked for the causes, then tried fixes, and kept the test months
untouched throughout.

\subsection{Why the training runs disagree on amounts}
\label{sec:why_runs}
A run's rain total can be wrong in three ways: it calls rain ``dry'' (missed rain), it calls a dry moment ``rain'' (false
rain), or it says ``rain'' at the right moments but with the wrong strength. We split the total of each of the five runs
of Section~\ref{sec:se_results} into these three parts, against the radar, on the July validation weeks and on the
August test month:
\begin{center}\footnotesize\setlength{\tabcolsep}{5pt}
\begin{tabular}{cccccc}
\toprule
& \multicolumn{2}{c}{\textbf{July (validation)}} & \multicolumn{3}{c}{\textbf{August (test)}}\\
\cmidrule(lr){2-3}\cmidrule(lr){4-6}
Run & bias & strength when it says rain & bias & strength & missed / false rain\\
\midrule
\SeedRows
\bottomrule
\end{tabular}
\end{center}
(``Strength'' is the run's rain divided by the radar's, at the moments when both say it rains; 1.00 is perfect. The
missed and false columns are shares of the radar's total.)

Three things stand out:
\begin{itemize}
\item \textbf{It is the strength, not the switch.} All runs miss about the same rain (10--12\%) and add about the same
false rain. What differs is how strongly they convert a decibel into millimetres: from 0.69 to 0.89 of the radar in August.
Fixing the rain/no-rain threshold at the same value for every run does not change this.
\item \textbf{July tells us about August.} A run that reads high in July also reads high in August: the correlation between
the two columns of bias is \SeedValTestR. So the spread is a property of each trained network, not bad luck in August.
\item \textbf{All runs read lower in August than in July}, by 15 to 20 points. August 2015 had more heavy downpours, and the
strength is learned mostly from the few heavy minutes of six weeks in one city. Each run fits those few minutes a little
differently.
\end{itemize}
That last point is the cause: too little and too uniform training rain. The cure we expected was more varied rain.

\subsection{What a long rain looks like to a link}
\label{sec:may6}
The Italian ``after a mostly rainy day'' hours are not spread over the summer. In 2022, 85\% of them come from a single
storm: a day and a half of steady rain on 5--7 May over the Apennine hills. Figure~\ref{fig:may6} shows one link through it,
next to a rain gauge 508\,m away.

\begin{figure}[h]
\centering
\begin{tikzpicture}
\begin{groupplot}[group style={group size=1 by 2,vertical sep=1.5cm,xlabels at=edge bottom},
  width=0.95\textwidth,height=4.6cm,xmin=0,xmax=60,grid=major,grid style={inkc!12},
  xtick={0,12,24,36,48,60},xticklabels={{5 May\\12:00},{6 May\\00:00},{6 May\\12:00},{7 May\\00:00},{7 May\\12:00},{8 May\\00:00}},
  xticklabel style={align=center,font=\scriptsize},
  tick label style={font=\small},label style={font=\small},
  legend style={font=\scriptsize,draw=none,fill=white,fill opacity=0.85,text opacity=1},legend cell align=left]
\nextgroupplot[ylabel={rain (mm/h)},ymin=0,ymax=4.5,legend pos=north east]
  \addplot[thick,color=gaugec,mark=*,mark size=1pt] table[x=hour,y=gauge] {may6.dat};
  \addlegendentry{rain gauge}
  \addplot[thick,color=radarc] table[x=hour,y=radar] {may6.dat};
  \addlegendentry{radar along the link}
  \addplot[thick,color=inkc,dashed] table[x=hour,y=ai_se] {may6.dat};
  \addlegendentry{AI, Sweden only}
  \addplot[thick,color=rainc] table[x=hour,y=ai_cv] {may6.dat};
  \addlegendentry{AI, varied training (never saw May)}
\nextgroupplot[ylabel={dB},xlabel={hourly averages (UTC)},ymin=0,ymax=6,legend pos=north east]
  \addplot[thick,color=inkc] table[x=hour,y=excess] {may6.dat};
  \addlegendentry{extra loss above the dry level}
  \addplot[thick,color=radarc,dashed] table[x=hour,y=expected] {may6.dat};
  \addlegendentry{loss the radar's rain should cause}
\end{groupplot}
\end{tikzpicture}
\caption{The long storm of 5--7 May 2022 on sub-link 182\_1 (5.3\,km, 24.6\,GHz). Top: rain. Bottom: what the link
measured. Storm totals: gauge \Maygauge\,mm, radar \Mayradar\,mm, AI trained on Sweden only \Mayaise\,mm, AI with varied
training \Mayaicv\,mm.}
\label{fig:may6}
\end{figure}

What the figure, and the same numbers for every link, show:
\begin{enumerate}
\item \textbf{The dry level is not the problem.} In every group of Italian rainy hours, the 3-day low value is at or a little below (by
0.1--0.2\,dB, and up to 0.9\,dB in the long storm) the level the link has when the radar says it is dry. It is never too
high, so it cannot be what makes the AI read too little.
\item \textbf{The link sees more than the rain.} While it rains, the extra loss is the rain's share plus a steady 1.5--2\,dB
more: water on the antenna covers, which stay wet in a long rain (the wet antenna of Section~\ref{sec:fakes}).
\item \textbf{The extra loss lasts a whole day after the radar stops} (6 May evening to 7 May noon). The gauge shows why:
light rain, 0.2--1.2\,mm/h, kept falling. The radar, looking from far above, saw almost none of it.
\item \textbf{The AI gives up.} After about twelve hours of raised loss, the AI reports almost nothing, even at hours when
the radar shows 2--4\,mm/h. To a single link, ``my loss has been raised for half a day'' looks like a faulty or drifting link, and
our training data are full of those. That is what the AI learned.
\end{enumerate}
So the cause is not one bug. Three things push the same way: a long rain wets the antennas, the radar we train on misses the
light rain at its end (so the training answer says ``dry'' while the link is still wet), and faulty links teach that a
loss which stays raised for hours is not rain.

\subsection{Trying fixes, one Italian summer at a time}
Each fix was trained five times (seeds 0--4) on Sweden June--July, plus Italy 2021 where stated, and tested on what it never
saw: Sweden in August and Italy in 2022. ``Rain after a mostly rainy day'' is measured here in 15-minute steps, which is
stricter than the hourly table of Section~\ref{sec:drylevel}.
\begin{center}\footnotesize\setlength{\tabcolsep}{3.5pt}
\begin{tabular}{p{5.1cm}ccccc}
\toprule
& \multicolumn{2}{c}{\textbf{Sweden, August}} & \multicolumn{3}{c}{\textbf{Italy 2022}}\\
\cmidrule(lr){2-3}\cmidrule(lr){4-6}
& & bias vs gauges, & total bias, & rain after a & rain after a\\
Training & MCC & 5 runs & 5 runs & dry day & mostly rainy day\\
\midrule
\ExpRows
\bottomrule
\end{tabular}
\end{center}
\begin{itemize}
\item \textbf{One Italian summer fixes most of the spread.} Five runs that disagreed by 47 points on Italy's 2022 total
now agree within 7 to 12 points, and in Sweden within 5 to 17 points instead of 26. Rain detection in Sweden improves too
(MCC 0.65 against 0.62), even though no Swedish data was added.
\item \textbf{Training only on reliable links helps the long storm.} With the reliability score of
Section~\ref{sec:reliability} applied to the Italian training links, the share of the long storm measured rises from 28\% to
42\%, with nothing else lost. This supports point 3 above: faulty links were teaching the AI to ignore long raised loss.
\item \textbf{Two ideas did not work.} Scoring the rate error in a log scale (so light rain counts as much as heavy rain)
made the AI read 5--15\% too little in Italy. Giving it two extra inputs from its neighbours (``how many links around me
are raised right now?'') improved rain detection (MCC 0.68) but cost amounts elsewhere and barely helped the long storm.
\end{itemize}

\subsection{The final check: every Italian month, judged by a model that never saw it}
\label{sec:cv}
One Italian summer may be luck. So we used both summers, month by month: to judge May, we train on Sweden plus June to
September of \emph{both} Italian years; to judge June, on the other four months; and so on. This is called
\textbf{cross-validation}. Every Italian month of 2021 and 2022 is then judged by a network that never saw it, and each
network learns from long rains of the other months (for example September 2021 when May is judged). Training uses the
reliable links only; the last week of each training month is kept apart to pick the checkpoint and the threshold; three
runs per month. Sweden's August is never used.

The same hourly table as in Section~\ref{sec:drylevel}, same links, same groups (rain measured as a share of the radar's):
\begin{center}\footnotesize
\begin{tabular}{lcccc}
\toprule
& \multicolumn{4}{c}{\textbf{share of the previous 24 hours that was rainy}}\\
\cmidrule(lr){2-5}
& under 10\% & 10--25\% & 25--50\% & over 50\%\\
\midrule
Italy 2021, Sweden-only AI & \SeOnlyLongAa & \SeOnlyLongAb & \SeOnlyLongAc & \SeOnlyLongAd\\
Italy 2021, \textbf{varied training} & \textbf{\CVrelLongAa} & \textbf{\CVrelLongAb} & \textbf{\CVrelLongAc} & \textbf{\CVrelLongAd}\\
\midrule
Italy 2022, Sweden-only AI & \SeOnlyLongBa & \SeOnlyLongBb & \SeOnlyLongBc & \SeOnlyLongBd\\
Italy 2022, \textbf{varied training} & \textbf{\CVrelLongBa} & \textbf{\CVrelLongBb} & \textbf{\CVrelLongBc} & \CVrelLongBd\\
\bottomrule
\end{tabular}
\end{center}
Against the rain gauges, over all rainy hours, the varied AI measures \CVrelGaugesA{} (2021) and \CVrelGaugesB{} (2022) of
the gauges' rain. The three runs per month agree on Italy's total within \CVrelBiasA{} (2021) and \CVrelBiasB{} (2022).
In Sweden's August, the same networks score MCC \CVrelSeMCC, hourly correlation with gauges \CVrelSeR, error
\CVrelSeRMSE\,mm/h and bias \CVrelSeBias{} (runs from \CVrelSeBiasRange), against 0.615, 0.79, 1.04\,mm/h and $-20\%$ for the
Sweden-only model: rain detection and agreement between runs improve clearly, the Swedish shortfall only a little.

\begin{careful}
\textbf{What this test is, and is not.} The Italian links are the same in every month: this tests new weather on known links,
not new links. The Italy and Netherlands results of Sections~\ref{sec:italy}--\ref{sec:neighbours}, and the maps of Part~3,
keep the Sweden-only model on purpose, so that they remain tests on countries the AI has never seen. For Algeria, which
has no reference data of its own, the recipe of this section is the one to use: train on several countries, reliable links
only, and average several runs.
\end{careful}

\begin{oneline}
More varied rain fixed most of both problems: training runs now agree on amounts, Sweden improved, and in Italy the AI
measures \CVrelLongAd{} of the rain after a rainy day in 2021. Two things remain: the one long storm of May 2022 is only
partly measured (\CVrelLongBd{}), and in Sweden's August the AI still reads \CVrelSeBias{} low.
\end{oneline}


%==========================================================
\section{The Huawei version: MindSpore}
\label{sec:mindspore}
%==========================================================

The competition asks for Huawei technology. Huawei's deep-learning toolkit is \textbf{MindSpore}; its AI chips are called
\textbf{Ascend}; and its cloud service for training models is \textbf{ModelArts}. We do not have Huawei Cloud access yet,
so the project is built so that the toolkit is one setting in a configuration file: \texttt{framework: torch} today,
\texttt{framework: mindspore} later. That setting already chooses which version of the network is built. Everything above
was trained with PyTorch, a widely used free toolkit.

To prove the switch works, we wrote the same network a second time in MindSpore, layer for layer, and ran two checks:
\begin{enumerate}
\item \textbf{Same answers.} We copied the 111{,}890 trained dials from the PyTorch network into the MindSpore one and
gave both the same three days of real signal from 64 Swedish links: 276{,}480 answers each. The largest difference in rain
was 0.00008\,mm/h, which is just the rounding of the computer's arithmetic; the rain/no-rain decision agreed on every single
minute.
\item \textbf{It learns.} We trained a fresh MindSpore network from random dials for 300 rounds on an ordinary processor
(no graphics card). Its penalty fell from 0.59 to 0.40: the MindSpore training loop works. What remains for a full run on Ascend
chips through ModelArts is to extend this short loop to the complete training recipe of Section~\ref{sec:tcn}
(validation checks, threshold choice), and cloud access.
\end{enumerate}

%==========================================================
\section{What Phase 2 taught us}
\label{sec:phase2_end}
%==========================================================

\begin{itemize}
\item \textbf{Rain detection: the AI is better in most tests} (Sweden, both Italian summers, Sweden with coarse
readings), without looking into the future. The Netherlands min/max test is the exception.
\item \textbf{Rain amounts: worse in Sweden, better abroad.} Trained on Sweden only, it measures about a fifth too little
in Sweden (the classical method 6\%), and this varies a lot between training runs. In Italy, where the classical method
breaks down, the AI's totals are within 10\% and its hourly correlation with gauges is 0.74--0.81.
\item \textbf{Varied training fixes most of the amount problem.} Trained also on Italian months it is never judged on,
the AI's training runs agree (Italian totals within a few percent of each other and of the radar), and Italian rain after a
rainy day is measured much better. With cleaned training answers (gauge-corrected radar, unclear moments hidden), Sweden's August shortfall shrinks
from $-20\%$ to \CVfinSeBias{} (Section~\ref{sec:varied}).
\item \textbf{Robustness: clearly better.} With coarse 1-dB readings, and in two countries it never saw, the AI made
no absurd readings in Italy and eight times fewer in the Netherlands. This matters for Algeria, where we will not be able
to retune anything by hand.
\item \textbf{Long rain: mostly fixed, one storm left.} A 3-day low dry level fixed much of the under-measurement after
rainy days in the Netherlands. In Italy the cause was elsewhere (wet antennas, a radar that misses light rain, faulty
training links); varied training on reliable links with cleaned answers fixed it in 2021 and mostly in 2022 (\CVfinLongBd), where one
day-long storm is still only partly measured. On Italian links the AI never saw, the recipe still helps, but less evenly (Section~\ref{sec:newlinks}).
\item \textbf{The physics conversion rescued our first designs}, which under-measured heavy rain badly. With the final
design, a network that outputs mm/h directly does about as well in Sweden; we have not tested it abroad.
\item \textbf{The typical link works across countries.} In Italy and the Netherlands, a typical link's hourly rain followed the
radar with correlation 0.76--0.84, close to Sweden's 0.88. The weak point is a minority of misbehaving links and hours.
\item \textbf{Neighbours help, link choice matters more.} Checking each link against its neighbours trims impossible hours
by a quarter to a half for both methods. Which links to trust matters more: a reliability score from each link's own raw
signal keeps two to three times more Italian links than hand-picking, with the AI scoring as well as on the hand-picked
set.
\end{itemize}

\textbf{Next, in order:}
\begin{enumerate}
\item Finding out why the AI still under-measures after rainy days in Italy, and why training runs differ so much in
amounts (more varied training data should help both). \emph{Done: Section~\ref{sec:varied}.}
\item Training on more varied data to improve the amounts. \emph{Done for Sweden + Italy (Sections~\ref{sec:cv}--\ref{sec:newlinks}); a test
leaving a whole country out still needs a third country with 1-minute data.}
\item Phase~3: combining many links into rain maps every few minutes, and pushing them forward in time for a short-range
forecast.
\end{enumerate}
